\practicechapter{第 12 章实践：控制流、跳转、循环、调用与 switch}

本章对应第一册第 12 章“控制流：跳转、循环、调用与 switch”。实践卷把控制流分成两层：机器码中的控制流，以及解析器源码中的错误路径控制流。本章重点训练后者，因为可靠工具首先要在坏输入面前走对路径。

\practicesection{一、对应原书问题：提前返回也是控制流设计}

控制流不只是汇编里的 \code{jmp}、\code{call}、\code{ret}。在解析器里，\code{inspect\_elf} 面对坏输入时必须选择是否继续。magic 不对，后续所有 ELF 字段都没有意义；header 太短，继续读取会越界；class 或 endianness 不支持，继续解释会得到错误字段。

因此本章把提前返回当成工程控制流来分析。每个提前返回都对应一个无法满足的前置条件，并且必须留下 diagnostic。没有 diagnostic 的提前返回会让用户不知道失败原因；没有提前返回的坏输入会污染后续证据。

\practicesection{二、从特例推到规律：坏输入路径要比好输入更清楚}

特例一：输入是普通文本文件。正确路径是设置 \code{valid\_magic=false}，追加 \code{input does not start with ELF magic}，返回。特例二：输入有 magic 但只有 16 字节。正确路径是承认 magic 正确，但追加 \code{input is shorter than ELF64 header}，返回。特例三：输入是 ELF32。正确路径是记录只支持 ELF64 小端，而不是尝试按 ELF64 读取。

由此推出规律：解析器的控制流应该由格式前置条件驱动。每通过一个前置条件，才能进入下一层读取。测试也要按这些条件组织，而不是只测试一个“正常文件能跑”。

\practicesection{三、实践任务：画出 \code{inspect\_elf} 的路径}

阅读 \code{inspect\_elf}，画出如下状态图的文字版：

\begin{enumerate}
  \item 创建 report，记录 source、byte size、checksum；
  \item 检查 magic，失败则 diagnostic 后返回；
  \item 检查 header 长度，失败则 diagnostic 后返回；
  \item 检查 class 和 endianness，失败则 diagnostic 后返回；
  \item 读取 header 字段；
  \item 解析 section；
  \item 解析 symbol；
  \item 返回完整 report。
\end{enumerate}

然后把测试文件中的三个坏输入测试与路径节点对应起来：\code{RejectsNonElfMagic} 对应第 2 步，\code{RejectsShortElfHeader} 对应第 3 步，\code{RejectsUnsupportedClassOrEndian} 对应第 4 步。最后在反汇编中观察一个真实函数调用，说明源码层调用、汇编层 \code{call} 和解析器控制流分析是不同层级的问题。

\practicesection{四、验收：错误路径不能含糊}

本章交付物是 \filepath{reports/ch12-control-flow.md}。通过标准如下：

\begin{enumerate}
  \item 能画出 \code{inspect\_elf} 的主要控制路径。
  \item 每个提前返回都写明触发条件和 diagnostic。
  \item 能解释为什么坏 magic 后不能继续读 header。
  \item 能说明测试覆盖了哪些错误路径，还没有覆盖哪些路径，例如 section table 越界。
  \item 能区分程序源码控制流、反汇编控制流和解析器错误处理控制流。
\end{enumerate}

这章的结论很实用：一个工具是否可靠，往往不是看正常输入有多快，而是看坏输入是否被清楚、稳定、可测试地处理。
